\Askhsh{32256}{4}
\begin{ylh}{1}
    \item 2.2 Παρουσίαση Στατιστικών Δεδομένων
    \item 2.3 Μέτρα Θέσης και Διασποράς
\end{ylh}
Στις τέσσερις κλάσεις του παρακάτω πίνακα, παρουσιάζονται τα χρόνια εργασίας του συνόλου των εκπαιδευτικών ενός Λυκείου της Αθήνας.

\begin{center}
\begin{mytblr}{}
{{Χρόνια\\Υπηρεσίας}} & {{Κεντρική}\\{Τιμή} $x_i$} & {{Συχνότητα $\nu_i$}} & $x_i \cdot \nu_i$ \\
$[0, 10)$ & & $10$ & \\
$[10, 20)$ & & $15$ & \\
$[20, 30)$ & & $κ$ & \\
$[30, 40)$ & & $5$ & \\

\textbf{Σύνολο} & \SetCell{bg=gray!20} & & \\
\end{mytblr}
\end{center}

\begin{alist}
\item Αφού μεταφέρετε στην κόλλα σας τον πίνακα, να συμπληρώσετε τα κενά του, συναρτήσει του $\kappa$, όπου αυτό είναι αναγκαίο. \monades{5}

\item Αν ο μέσος χρόνος εργασίας των παραπάνω εκπαιδευτικών είναι τα $19$ χρόνια, να δείξετε ότι το $\kappa=20$. Πόσοι ήταν συνολικά οι εκπαιδευτικοί του σχολείου; \monades{6}
\end{alist}
Για $\kappa=20$,
\begin{alist}\setcounter{enumi}{2}
\item να μεταφέρετε στην κόλλα σας τον παρακάτω πίνακα και να συμπληρώσετε τα κενά. \monades{8}

\begin{center}
\begin{mytblr}{}
{Χρόνια\\Υπηρεσίας} & $\nu_i$ & $f_i$ & $N_i$ & $F_i$ & $F_i\%$ \\
$[0, 10)$ & $10$ & & & & \\
$[10, 20)$ & $15$ & & & & \\
$[20, 30)$ & $20$ & & & & \\
$[30, 40)$ & $5$ & & & & \\

\textbf{Σύνολο} & &  & \SetCell{bg=gray!20} & \SetCell{bg=gray!20} & \SetCell{bg=gray!20} \\
\end{mytblr}
\end{center}

\item να κατασκευάσετε το ιστόγραμμα και το πολύγωνο των αθροιστικών σχετικών συχνοτήτων επί τοις εκατό, και να βρείτε τη διάμεσο του χρόνου εργασίας. \monades{6}
\end{alist}

